% STATUS: DRAFT
\chapter{Topologies on operators}\label{ch:topologies}

\begin{physmotiv}
When we say ``the observable $A$ is the limit of local observables $A_N$'' we must say in \emph{what sense}. The average magnetization $M_N=\frac1N\sum_{j=1}^N\sigma^z_j$ is a perfectly good local observable for each $N$; in a state with magnetization $m$ it converges to the number $m$ in an appropriate sense, but not in operator norm. This difference between the modes of convergence is exactly the difference between a \cstar-algebra (closed in norm) and a von Neumann algebra (closed in weak/strong topologies). Much of the structure of the book depends on this chapter.
\end{physmotiv}

\section{Three ways an operator sequence can converge}

Let $A_n,A\in\BH$.
\begin{center}
\renewcommand{\arraystretch}{1.3}
\begin{tabular}{@{}lll@{}}
\toprule
Mode & Definition & Physical meaning\\
\midrule
\textbf{norm} & $\norm{A_n-A}\to0$ & uniform precision on all states\\
\textbf{strong} & $\norm{(A_n-A)\xi}\to0\;\forall\xi\in\HH$ & each state is mapped to a convergent state\\
\textbf{weak} & $\inner\eta{A_n\xi}\to\inner\eta{A\xi}\;\forall\xi,\eta$ & all matrix elements/expectation values converge\\
\bottomrule
\end{tabular}
\end{center}
Norm convergence implies strong convergence implies weak convergence, and the converses fail.

\begin{example}[the three modes are different]\label{ex:shift}
Let $P_n$ be the projection onto $\mathrm{span}(e_0,\dots,e_{n-1})$ in $\ell^2$. Then $P_n\to\id$ strongly (since $\norm{(\id-P_n)\xi}^2=\sum_{k\ge n}|\xi_k|^2\to0$), but $\norm{\id-P_n}=1$ for all $n$, so $P_n\not\to\id$ in norm. Next, the shift powers $(S\ad)^n\to0$ strongly (the tail of a vector goes to zero) but not in norm. And $S^n\to0$ \emph{weakly} ($\inner\eta{S^n\xi}\to0$ since the shifted vector escapes to infinity) but \emph{not} strongly, since $\norm{S^n\xi}=\norm\xi$.
\end{example}

\begin{whatwrong}
Multiplication $(A,B)\mapsto AB$ is norm-continuous, but only \emph{separately} strongly continuous (and only jointly so on bounded sets); for the weak topology it is only separately continuous. In \cref{ex:shift}, $S^n\to0$ weakly and $(S\ad)^n\to0$ weakly yet $(S\ad)^nS^n=\id$ for all $n$. Never take a product of weak limits without checking.
\end{whatwrong}

For general (non-separable or non-sequential) settings one needs \emph{nets} instead of sequences; for bounded sets in separable $\HH$ sequences suffice. A basis of neighbourhoods of $A$ in the strong topology is given by sets $\{B:\norm{(B-A)\xi_i}<\epsilon,\ i=1,\dots,k\}$ for finitely many vectors, and in the weak topology by $\{B:|\inner{\eta_i}{(B-A)\xi_i}|<\epsilon\}$.

\section{A physics example: macroscopic observables become classical}

Consider the infinite spin chain in the product state $\Psi_\theta$ from \cref{ch:why_algebras} (with magnetization $m=\cos2\theta$). Let $M_N=\frac1N\sum_{j=1}^N\sigma^z_j$.

\begin{itemize}
\item $\norm{M_N}=1$ for all $N$ in the $N$-site algebra. It does \emph{not} converge in norm: $\norm{M_N-M_{N'}}$ is of order 1 for $N'\gg N$ (choose a state aligned with the first $N$ sites and orthogonal on the remainder).
\item In the representation built on $\Psi_\theta$ (the GNS representation, \cref{ch:gns}), the variance is $\norm{(M_N-m)\Psi_\theta}^2=\frac{1-m^2}{N}\to0$. Since the local vectors $A\Psi_\theta$ are dense and $M_N$ commutes with all but finitely many sites, $M_N\to m\,\id$ \emph{strongly}.
\end{itemize}
So in the representation $\pi_\theta$, the limit of $M_N$ exists and is a \emph{scalar} $m\id$, which commutes with every operator. Different values of $m$ give different representations, so $M=\lim M_N$ is a ``classical'' variable that labels the sector. In algebraic language: \emph{macroscopic observables become central elements in the von Neumann algebra obtained by weak closure}. The set of such central elements (the centre of the algebra) is the algebraic encoding of ``which phase are we in?'' (\cref{ch:factors}).

\begin{keyidea}
Norm closure (\cstar-algebra) keeps the quantum structure of local observables; weak/strong closure (von Neumann algebra) adds macroscopic/classical observables as limits, and these live in the centre.
\end{keyidea}

\section{The weak topologies and states}

\begin{definition}
The \emph{$\sigma$-weak} (ultraweak) topology on $\BH$ is the weak-$^*$ topology it inherits as the dual of $\mathcal T(\HH)$: $A_\alpha\to A$ iff $\Tr(\rho A_\alpha)\to\Tr(\rho A)$ for every density matrix $\rho$. The \emph{$\sigma$-strong} topology is defined similarly using sequences of vectors with $\sum\norm{\xi_k}^2<\infty$.
\end{definition}
On bounded sets the weak and $\sigma$-weak topologies coincide, and the same for strong and $\sigma$-strong. A linear functional on $\BH$ is $\sigma$-weakly continuous iff it is of the form $A\mapsto\Tr(\rho A)$ for a trace-class $\rho$. In particular:

\begin{center}
\emph{normal states $=$ $\sigma$-weakly continuous states $=$ states given by density matrices.}
\end{center}
This is what ``normal'' will mean from now on (\cref{ch:traces}); the singular states of \cref{ch:hilbert} are the ones that fail it.

\begin{theorem}[Banach--Alaoglu in this setting]
The unit ball of $\BH$ is compact in the weak (and $\sigma$-weak) topology.
\end{theorem}
Compactness of the unit ball is the reason limits of bounded sequences of observables have weak limit points. Because weak limits of normalized states can be non-normal, one needs to be careful about which topology one uses for states vs observables.

\section{Why this matters}

\begin{itemize}
\item The \emph{double commutant theorem} (\cref{ch:doublecomm}) states that for a unital $^*$-algebra $\M\subset\BH$ the weak, strong and $\sigma$-weak closures coincide with $\M''$. Von Neumann algebras are defined by this property, equivalently by being closed in any of these topologies.
\item Normality (continuity of a state in the $\sigma$-weak topology) is the correct notion of ``physical'' state for a vN algebra.
\item The Hilbert space chosen matters for \emph{closure}: the weak closure of the same \cstar-algebra in different representations gives different vN algebras (e.g.\ product states with different $m$).
\end{itemize}

\section*{Exercises}
\begin{exercise}
Verify the claims of \cref{ex:shift}: $P_n\to\id$ strongly but not in norm; $S^n\to0$ weakly but not strongly; $(S\ad)^n\to0$ strongly but not in norm.
\end{exercise}
\begin{exercise}
Show that for the product state $\Psi_\theta$ one has $\norm{(M_N-m)\Psi_\theta}^2=(1-m^2)/N$. Deduce that $M_N\to m\id$ strongly in the representation on the space generated from $\Psi_\theta$.
\end{exercise}
\begin{exercise}
In the same setting, show that $M_N$ is not Cauchy in norm. (Hint: compare $M_N$ and $M_{2N}$ on a state with the first $N$ spins up and the next $N$ down.)
\end{exercise}
\begin{exercise}
Let $f_n(x)=\ee^{\ii nx}$ acting by multiplication on $L^2([0,2\pi])$. Does $f_n\to0$ weakly? Strongly? In norm? (Riemann--Lebesgue.)
\end{exercise}
\begin{exercise}
Explain why the product of weakly convergent sequences need not converge to the product of the limits, using $S^n$ and $(S\ad)^n$.
\end{exercise}
